\begin{table}[t]
\centering

\begin{tabular}{lccc}
\toprule
\textbf{Method} & \textbf{Controllability} & \textbf{Multiplayer} & \textbf{Graphics} \\
\midrule
Malmo~\cite{johnson2016malmo} & \texttimes & \checkmark & \checkmark \\
MineRL~\cite{guss2019minerllargescaledatasetminecraft}     & \texttimes & \texttimes & \checkmark \\
MineDojo~\cite{fan2022minedojo}  & \texttimes & \texttimes & \checkmark \\
Voyager~\cite{wang2023voyager} & \checkmark & \texttimes & \texttimes \\
Mineflayer~\cite{mineflayer2025} & \checkmark & \checkmark & \texttimes \\
\solarisengine    & \checkmark & \checkmark & \checkmark \\
\bottomrule
\end{tabular}
\caption{\textbf{Comparison of Minecraft AI frameworks.} Unlike prior systems, \solarisengine enables controlled multiplayer gameplay collection with real Minecraft graphics. RL-based frameworks such as Malmo, MineRL, and MineDojo produce visual observations but offer limited controllability, as their low-level action spaces require training RL agents to collect meaningful data, which is prohibitively expensive and counterproductive for world-modeling data collection. Voyager and Mineflayer provide high-level behavioral control but operate in text-only mode without visual output. \solarisengine supports multiplayer gameplay, fine-grained programmatic control and rich visual observations.}
\label{tab:mc_systems_comparison}
\end{table}
